%% kapa: exam=2 round=35 year=2024 period=1 subject=appraisal_practice subject_ko=감정평가실무 seq=2 number=2 points=30
%% source: Q-Net 감정평가사 기출문제 (cst00309) artlSeq=5213532 file=raw/1008/5213532/2024년 제35회 감정평가사 2차시험 문제지.pdf
%% conversion: kordoc 4.24.51 -> tools/gen2.py (auto); review=proofread-sonnet; flags=none
\begin{longquestion}{【문제 2】}
상업용부동산의 감정평가에 대하여 주어진 자료를 바탕으로 물음에 답하시오. (30점) \par
\medskip\noindent 물음 1) 시산가액 조정기준(reconciliation criteria)을 설명하고, 이와 관련지어 본건 시산가액 조정을 통한 감정평가액 결정의 적정성에 대한 의견을 기술하시오. (10점) \par
\medskip\noindent 물음 2) 주어진 자료 상 본건에 적용한 환원율이 부적정할 가능성이 있는 사유에 대하여 설명하고, 제시된 자료 외에 추가적으로 부적정한 원인이 될 수 있는 사유를 기술하시오. (10점) \par
\medskip\noindent 물음 3) 인근지역의 시장상황에 비추어 적정하다고 판단되는 환원율을 산출하고, 이를 기초로 산출된 수익방식에 의한 시산가액이 대상부동산의 적정한 감정평가액(시장가치)일 경우 본건 건물의 유효잔존내용연수를 구하시오. (내용연수 만료 시 잔존가치와 철거비는 없음) (10점) \par
\medskip\noindent < 자료 > 감정평가 요약 \par
1. 기본적 사항 \par
\hspace*{0.5em}- 소 재 지: S시 S구 B동 157 \par
\hspace*{0.5em}- 토\hspace{1.2em}지: 3종일반주거지역, 대, 500m\textsuperscript{2} \par
\hspace*{0.5em}- 건\hspace{1.2em}물: 연와조 슬래브지붕, 지하1층 지상2층, 연면적 900m\textsuperscript{2}, 사용승인일 1994. 07. 13. \par
\hspace*{0.5em}- 평가목적: 일반거래(시가참고) \par
\hspace*{0.5em}- 기준시점: 2024. 07. 13. \par
\hspace*{0.5em}- 기준가치: 시장가치 \par
2. 인근지역 및 본건 개황 \par
\hspace*{0.5em}1) 인근지역 개황 \par
\hspace*{1em}- 본건의 인근지역은 S시 부도심 서측의 N로(중로)에 접하고 있는 노선상가지대로, 최근 개통된 터널로 S시 부도심의 간선도로와 동서로 연결되었음 \par
\hspace*{1em}- 인근지역의 부동산은 그동안 낮은 상업성으로 인해 저밀도로 이용되고 있었는데, 터널의 개통과 함께 경과연수가 오래된 건물들부터 철거 후 신축 또는 대수선을 수반한 증축이 일어나고 있어 지역의 변모와 함께 부동산 거래량이 증가하는 양상을 나타내고 있음 \par
\hspace*{1em}- 인근지역 시장조사 결과, 유사부동산에 대한 투자수익률은 8.00\%이고 순임료(NOI)와 부동산가격이 매년 3.00\% 상승할 것으로 예상됨 \par
\hspace*{0.5em}2) 대상 부동산 개황 \par
\hspace*{1em}- 본건 토지는 N로에 북측으로 접하고 있는 가장형 평지임 \par
\hspace*{1em}- 본건 건물은 30년 경과된 지하1층 지상2층 연와조 슬래브지붕으로, 신축 후 양호한 관리 및 소매판매점 용도에 비추어 시간의 경과에 따른 감가 외에 기능적 감가는 없는 것으로 판단됨 \par
3. 원가방식에 의한 시산가액 산출 \par
\hspace*{0.5em}1) 토지가치의 산출 \par
\hspace*{1em}- 공시지가기준법에 의한 토지가치: @18,000,000원/m\textsuperscript{2} \par
\hspace*{1em}- 거래사례비교법에 의한 토지가치: @18,200,000원/m\textsuperscript{2} \par
\hspace*{1em}- 거래사례비교법에 의하여 충분히 지지되는 공시지가기준법에 의한 토지가치로 결정 \par
\hspace*{1em}- @18,000,000원/m\textsuperscript{2} $\times$ 500m\textsuperscript{2} = 9,000,000,000원 \par
\hspace*{0.5em}2) 건물가치의 산출 \par
\hspace*{1em}- 재조달원가: @1,500,000원/m\textsuperscript{2}(표준건축비에 따른 적정한 수준) \par
\hspace*{1em}- 감가수정(정액법): @1,500,000원/m\textsuperscript{2} $\times$ 30년 / 45년 = @1,000,000원/m\textsuperscript{2} \par
\hspace*{1em}- 건물가치: (@1,500,000원/m\textsuperscript{2} - @1,000,000원/m\textsuperscript{2}) $\times$ 900m\textsuperscript{2} = 450,000,000원 \par
\hspace*{0.5em}3) 원가방식에 의한 시산가액 \par
\hspace*{1em}토지가치: 9,000,000,000원 \par
\hspace*{1em}건물가치:\hspace{0.8em}450,000,000원 \par
\hspace*{1em}합\hspace{1em}계: 9,450,000,000원 \par
4. 수익방식에 의한 시산가액 산출 \par
\hspace*{0.5em}1) 계약임대료: @42,500원/m\textsuperscript{2}, 월(순임료) \par
\hspace*{0.5em}2) 계약내용: 본건 부동산 전체를 할인마트에 장기임대 중으로 최근에 시장임대료 수준으로 계약임대료를 갱신하였는바, 잔존 임대차기간은 7년임 \par
\hspace*{0.5em}3) 순영업소득: @42,500원/m\textsuperscript{2} $\times$ 900m\textsuperscript{2} $\times$ 12 = 459,000,000원 \par
\hspace*{0.5em}4) 환원율의 산출 \par
{\small\centering
\begin{tabular}{|c|c|c|c|c|c|}
\hline
\makecell{매매\\사례} & 매매가격(원) & \makecell{순영업소득(원)\\(NOI)} & \makecell{환원율\\(Ro)} & 토지가격(원) & 건물가격(원) \\
\hline
1 & 8,500,000,000 & 488,000,000 & 5.74\% & 5,400,000,000 & 3,100,000,000 \\
\hline
2 & 5,300,000,000 & 300,000,000 & 5.66\% & 3,200,000,000 & 2,100,000,000 \\
\hline
3 & 14,000,000,000 & 798,000,000 & 5.70\% & 8,700,000,000 & 5,300,000,000 \\
\hline
\end{tabular}\par}

\hspace*{0.5em}※ 상기 매매사례는 모두 3종일반주거지역에 속하고 매매가격, 토지 및 건물가격, 순영업소득 등은 모두 정상적인 최근 사례임 \par
\hspace*{0.5em}최근 본건 인근에 소재하는 상기 매매사례의 매매가격과 순영업소득을 기초로 산출된 환원율이 유사한 수준을 나타내고 있어 신뢰성이 있다고 판단되어 본건 평가에 적용할 환원율을 5.70\%로 결정함 \par
\hspace*{0.5em}5) 수익방식에 의한 시산가액 \par
\hspace*{1em}459,000,000원 / 5.70\% $\doteqdot$ 8,052,000,000원 \par
5. 시산가액의 조정 및 감정평가액 결정 \par
\hspace*{0.5em}- 원가방식에 의한 시산가액은 거래사례비교법에 의하여 충분히 지지되는 공시지가기준법에 의하여 산출된 시산가액에 정액법에 의하여 감가수정한 건물가치를 합산하여 산출하였음 \par
\hspace*{0.5em}- 수익방식에 의한 시산가액은 최근 시장임대료 수준으로 갱신된 계약임대료를 기초로 산출된 순영업소득에 인근의 정상적인 매매사례의 순영업소득과 매매가격에 기초하여 산출된 환원율 5.66\% \textasciitilde{} 5.74\%를 고려하여 환원율 5.70\%를 적용하여 산출하였음 \par
\hspace*{0.5em}- 이와 같이 원가방식 및 수익방식에 적용한 자료 등이 모두 기준시점 현재 인근의 시장상황을 반영하고 있는바, 각 감정평가방식에 특별한 문제가 없는 것으로 판단됨 \par
\hspace*{0.5em}- 따라서 양 시산가액의 평균액을 최종 감정평가액 8,751,000,000원으로 결정함 \par
\end{longquestion}
